%-------------------------
% Resume in Latex
% Author : Brian Pulfer (Template by Jake Gutierrez)
% Based off of: https://github.com/sb2nov/resume
% License : MIT
%------------------------

\documentclass[letterpaper,11pt]{article}

\usepackage{latexsym}
\usepackage[empty]{fullpage}
\usepackage{titlesec}
\usepackage{marvosym}
\usepackage[usenames,dvipsnames]{color}
\usepackage{verbatim}
\usepackage{enumitem}
\usepackage[hidelinks]{hyperref}
\usepackage{fancyhdr}
\usepackage[english]{babel}
\usepackage{tabularx}
\usepackage{graphicx}
\usepackage{float}
\input{glyphtounicode}


%----------FONT OPTIONS----------
% sans-serif
% \usepackage[sfdefault]{FiraSans}
% \usepackage[sfdefault]{roboto}
% \usepackage[sfdefault]{noto-sans}
% \usepackage[default]{sourcesanspro}

% serif
% \usepackage{CormorantGaramond}
% \usepackage{charter}


\pagestyle{fancy}
\fancyhf{} % clear all header and footer fields
\fancyfoot{}
\renewcommand{\headrulewidth}{0pt}
\renewcommand{\footrulewidth}{0pt}

% Adjust margins
\addtolength{\oddsidemargin}{-0.5in}
\addtolength{\evensidemargin}{-0.5in}
\addtolength{\textwidth}{1in}
\addtolength{\topmargin}{-.5in}
\addtolength{\textheight}{1.0in}

\urlstyle{same}

\raggedbottom
\raggedright
\setlength{\tabcolsep}{0in}

% Sections formatting
\titleformat{\section}{
  \vspace{-4pt}\scshape\raggedright\large
}{}{0em}{}[\color{black}\titlerule \vspace{-5pt}]

% Ensure that generate pdf is machine readable/ATS parsable
\pdfgentounicode=1

%-------------------------
% Custom commands
\newcommand{\resumeItem}[1]{
  \item\small{
    {#1 \vspace{-2pt}}
  }
}

\newcommand{\resumeSubheading}[4]{
  \vspace{-2pt}\item
    \begin{tabular*}{0.97\textwidth}[t]{l@{\extracolsep{\fill}}r}
      \textbf{#1} & #2 \\
      \textit{\small#3} & \textit{\small #4} \\
    \end{tabular*}\vspace{-7pt}
}

\newcommand{\resumeSubheadingNoPlace}[2]{
  \vspace{-2pt}\item
    \begin{tabular*}{0.97\textwidth}[t]{l@{\extracolsep{\fill}}r}
      \textbf{#1} & #2 \\
    \end{tabular*}\vspace{-7pt}
}

\newcommand{\resumeSubSubheading}[2]{
    \item
    \begin{tabular*}{0.97\textwidth}{l@{\extracolsep{\fill}}r}
      \textit{\small#1} & \textit{\small #2} \\
    \end{tabular*}\vspace{-7pt}
}

\newcommand{\resumeProjectHeading}[2]{
    \item
    \begin{tabular*}{0.97\textwidth}{l@{\extracolsep{\fill}}r}
      \small#1 & #2 \\
    \end{tabular*}\vspace{-7pt}
}

\newcommand{\resumeSubItem}[1]{\resumeItem{#1}\vspace{-4pt}}

\renewcommand\labelitemii{$\vcenter{\hbox{\tiny$\bullet$}}$}

\newcommand{\resumeSubHeadingListStart}{\begin{itemize}[leftmargin=0.15in, label={}]}
\newcommand{\resumeSubHeadingListEnd}{\end{itemize}}
\newcommand{\resumeItemListStart}{\begin{itemize}}
\newcommand{\resumeItemListEnd}{\end{itemize}\vspace{-5pt}}


%-------------------------------------------
%%%%%%  RESUME STARTS HERE  %%%%%%%%%%%%%%%%%%%%%%%%%%%%

\begin{document}
%----------HEADING----------

%\begin{center}
%\textbf{\Huge \scshape Brian Pulfer} \\ \vspace{1pt}
%\end{center}
\begin{center}
    \textbf{\Huge \scshape Brian Pulfer} \\ \vspace{1pt}
    \href{mailto:me@brianpulfer.ch}{\underline{me@brianpulfer.ch}} \\ \vspace{1pt}
    \href{https://www.brianpulfer.ch}{\underline{brianpulfer.ch}} $|$
    \href{https://brianpulfer.ch/github}{\underline{brianpulfer.ch/github}} $|$
    \href{https://brianpulfer.ch/linkedin}{\underline{brianpulfer.ch/linkedin}} $|$
    \href{https://brianpulfer.ch/scholar}{\underline{brianpulfer.ch/scholar}}
   \\ \vspace{1pt}
    \small{Work Authorization: U.S. O-1 visa}
   \\
\end{center}

%-----------EXPERIENCE-----------
\section{Experience}
  \resumeSubHeadingListStart

    \resumeSubheading
      {Research Scientist, Instagram Reels Translation}{Oct. 2026 -- Present}
      {Meta Platforms, Inc.}{Menlo Park, California}

    \resumeSubheading
      {Research Scientist Intern}{June 2025 -- September 2025}
      {Meta Platforms, Inc.}{Menlo Park, California}
      \resumeItemListStart
      \resumeItem{NVFP4 quantization of LLMs on the vLLM inference engine, cutting memory by \textbf{67.5\%}/\textbf{37.6\%} (vs.\ BF16/FP8) with $< 2\%$ accuracy drop.}
      \resumeItem{Exploratory studies on NVFP4 quantization of KV-cache in vLLM.}
      \resumeItemListEnd
  
    \resumeSubheading
      {Graduate Researcher (Ph.D. in Machine Learning)}{November 2021 -- September 2026}
      {University of Geneva}{Geneva, Switzerland}
      \resumeItemListStart
      \resumeItem{Dissertation on self-supervised vision models and robustness to adversarial attacks; published at CVPRW 2025 and ECCV 2024.}
      \resumeItem{Teaching Assistant for Algorithms and Data Structures courses.}
      \resumeItemListEnd
\resumeSubHeadingListEnd


%-----------EDUCATION-----------
\section{Education}
  \resumeSubHeadingListStart
    \resumeSubheading
      {UNIGE - University of Geneva}{Geneva, Switzerland}
      {Ph.D. in Machine Learning}{Nov. 2021 -- Sep. 2026}
    \resumeSubheading
      {USI - University of Southern Switzerland}{Lugano, Switzerland}
      {MSc in Artificial Intelligence (GPA 9.1/10)}{Aug. 2019 -- Jun 2021}
    \resumeSubheading
      {SUPSI - University of Applied Sciences of Southern Switzerland}{Manno, Switzerland}
      {BSc in Computer Science}{Aug. 2016 -- Jun 2019}
  \resumeSubHeadingListEnd


%-----------Publications-----------
\section{Selected Publications \small{[Full list: \href{https://brianpulfer.ch/scholar}{Google Scholar}]}}
    \resumeItemListStart
        \resumeItem{\textbf{B. Pulfer} et al. \textit{Task-Agnostic Attacks Against Vision Foundation Models.} CVPRW 2025 (AdvML Workshop).}
    
        \resumeItem{\textbf{B. Pulfer} et al. \textit{Robustness Tokens: Towards Adversarial Robustness of Transformers.} ECCV 2024.}
        
        \resumeItem{A. Stocco, \textbf{B. Pulfer}, P. Tonella. \textit{Mind the Gap! Transferability of Virtual vs.\ Physical-world Testing of Autonomous Driving Systems.} IEEE TSE.}

    \resumeItemListEnd

%-----------HONORS & SERVICE-----------
\section{Honors \& Service}
  \begin{itemize}[leftmargin=0.15in, label={}]
    \small{\item{
      Reviewer for CVPR 2025; Judge at the \href{https://neuralwave.ch}{Neural Wave} hackathon; SODESKA merit scholarship (USI, 2021); 3rd place, NeurIPS Weather4cast Challenge 2022; Winner, HackZ\"urich 2023.
    }}
  \end{itemize}


%-----------OPTIONAL-----------
\begin{comment}
\section{Others}
    \resumeItemListStart
        \resumeItem{Moderator of \textit{The AI Epiphany}'s community discord server}
        
    \resumeItemListEnd
\end{comment}



%-----------TECH SKILLS-----------
\section{Technical Skills}
 \begin{itemize}[leftmargin=0.15in, label={}]
    \small{\item{
     \textbf{Languages}{: Python, C++/CUDA} \\
     \textbf{Frameworks}{: PyTorch, vLLM, Hugging Face} \\
     \textbf{Areas}{: LLM inference \& quantization, machine translation, adversarial robustness, distributed training}
    }}
 \end{itemize}

%-----------PROJECTS-----------
\begin{comment}
\section{Projects}
    \resumeSubHeadingListStart
      \resumeProjectHeading
          {\textbf{Master Thesis} $|$ \emph{Python3, C\#, Keras, pandas, Unity, Git, Conda, cv2}}{September 2020 -- June 2021}
          \resumeItemListStart
            \resumeItem{Assembled a physical DonkeyCar using a JetsonNano computer and an RC car.}
            \resumeItem{Created a simulated scene of a real-world lab room in Unity. Improved the Unity simulator to log testing metrics.}
            \resumeItem{Collected data, trained and tested different DL models for self-driving in the simulated and real world tracks.}
            \resumeItem{Adapted CycleGAN to translate simulated images to real ones and train a real-world Cross-track-error predictor.}
          \resumeItemListEnd
      \resumeProjectHeading
          {\textbf{Bachelor Thesis} $|$ \emph{Python3, Git, Unittest, Gensim, PyJNIus, Keras, SkLearn}}{May 2019 -- Sep 2019}
          \resumeItemListStart
            \resumeItem{Developed a binary classifier machine learning model that can tell if two scientific articles from the PUBMED database were published by the same author. Work commissioned by La Roche AG.}
            \resumeItem{Implemented feature extraction code, also using a Java library inpython through the PyJNIus library. Used the Gensim library to apply doc2vec techniques, a novelty in the literature of AND.}
            \resumeItem{Trained and tested different models: KNN, SVM, Random Forest and Feed-Forward Neural Networks.}
            \resumeItem{Studied ambiguity level in the PubMed dataset by counting the cardinality of the namespaces.}
          \resumeItemListEnd
    \resumeSubHeadingListEnd
\end{comment}
 
%\bibliographystyle{plain}
%\bibliography{refs}
\end{document}
